\subsection{Rule Retrieval}
\label{sec:retrieval}

At inference time, injecting all learned rules is inefficient and noisy, making retrieval
necessary to select only rules relevant to the current query and tool context.


\subsubsection{Natural language based}

A simple baseline retrieves rules directly from their natural-language form by encoding the
query and available tools with an off-the-shelf sentence embedding model and selecting the
top-$k$ rules by cosine similarity. While easy to implement, this approach relies on generic
embeddings that struggle with lexical variability and fine-grained rule semantics. Finetuning
a dedicated embedding model introduces additional complexity, including limited training
data, the need for ground-truth rule annotations, and reduced scalability.

\subsubsection{Symbolic-based}

To address these limitations, and inspired by prior work on structure-aware retrieval
\citep{munikoti2023atlantic,li2023santa}, we perform retrieval over the symbolic
representation of rules. By structuring both rules and queries into shared semantic fields,
symbolic matching is more efficient, interpretable, and more robust to incidental linguistic
variation than retrieval over raw text.

Given a query \( x_q \) and available tools \( \mathcal{S}_q \), we first convert the state into a symbolic representation \( z_q \). Retrieval proceeds in two stages. First, a coarse filter removes inapplicable rules: decomposition rules are always retained, while tool-use rules are filtered based on whether their scope matches the available tools or tool categories. Second, the remaining rules are ranked by semantic similarity between their symbolic fields and those of \( z_q \), using a generic embedding model. We retain the top-$k$ rules and inject them into the prompt in natural-language form. Unlike retrieval over raw text, similarity is computed over a fixed symbolic vocabulary, where NL-to-symbolic conversion has already resolved much of the lexical and structural variation that generic embeddings struggle to capture.